凸性保证了走不动的地方就是最优点，但它不负责把你送到那里。从一个初值到一个可信的解，中间是几百次乃至几百万次迭代，每一次都要回答同一个问题：站在当前点，我掌握的关于这个函数的信息，够支持我往哪个方向、走多远。

五种方法的差别几乎全在这句话里。梯度下降只拿到一条切线，方向免费而长度昂贵，能走多远由曲率上界写死；Newton 法多拿到一个 Hessian，方向被曲率校正，步长这个旋钮随之消失；目标带尖角时连梯度都不存在，次梯度把切线换成一族支撑直线，每一步于是不再保证下降；若那个尖角的形状是已知的，近端算子可以把它精确求解而不是粗糙近似，速率随之回到光滑情形的水平；约束把可行域围成一块难以投影的区域时，障碍函数把边界改写进目标，让迭代自己待在里面。

这条线索有一个统一的读法：算法的强弱不取决于更新公式写得多复杂，而取决于它调用了目标函数的哪个接口。只能查询次梯度，就只能走 \(O(1/\sqrt k)\)；能查询梯度并知道光滑常数，就有 \(O(1/k)\)；能形成并分解 Hessian，就有局部二次收敛。接口越强，单步越贵，收敛越快——这三件事是绑在一起的，所以``哪个优化器最好''这个问题在脱离问题结构时没有答案。

\subsection{怎么读这一章}

按信息量递增的顺序读。第一篇建立整章的度量标尺：光滑常数 \(L\) 给出步长上界，条件数 \(\kappa\) 给出迭代次数，后面每一种方法都是在改善这两个数中的一个。第二篇给出改善条件数的极端做法，顺便交代一个贯穿全章的观念——对坐标缩放不变的算法不需要调步长，这也是深度学习里自适应方法在模仿的东西。第三篇和第四篇处理同一类非光滑目标，区别只在于把不可微项当黑箱还是当已知结构，两者的速率差整整一档，值得对照着读。第五篇回到带约束的一般情形，并把\hyperref[ch:075]{``对偶：为最优性建立下界证书''}里的对偶间隙变成一个可以提前规划的停止尺度。

每读完一篇，问三个问题：它向目标函数索取了什么，一步的最大开销花在哪一行，以及停下来时能交出什么样的证据。这三个答案凑齐了，选型就不再靠名字。

\subsection{本章篇目}

以下篇目按概念依赖排列。

\begin{itemize}
\tightlist
\item \hyperref[sec:09-Convex-Optimization/06-Optimization-Algorithms/01]{``梯度下降真正难在步长''}；
\item \hyperref[sec:09-Convex-Optimization/06-Optimization-Algorithms/02]{``Newton 法用曲率校正方向''}；
\item \hyperref[sec:09-Convex-Optimization/06-Optimization-Algorithms/03]{``不可微时次梯度为何仍能前进''}；
\item \hyperref[sec:09-Convex-Optimization/06-Optimization-Algorithms/04]{``近端与坐标方法利用可分结构''}；
\item \hyperref[sec:09-Convex-Optimization/06-Optimization-Algorithms/05]{``从障碍函数到原始---对偶内点法''}。
\end{itemize}

前四篇解释一个学习率为什么值得调一整天，最后一篇解释商用求解器为什么一个参数都不用你调。
